\section{实验结论}
本次实验成功实现了两个任务的预定目标:

\textbf{任务一}完成了基于内存的类FAT-16文件系统的设计与实现。该系统以数组\texttt{Disk[DISK\_MAXLEN]}模拟硬盘空间,在此基础上定义了FAT表、目录项(DirEntry)和文件描述符(FileDescriptor)等核心数据结构,并实现了目录创建(mkdir)、目录查看(ls)、目录切换(cd)、文件创建(touch)、文件打开(open)、文件读取(read)、文件写入(write)、文件关闭(close)以及文件删除(rm)等完整的基本文件系统操作。经测试验证,各项功能均能正常运行,表明FAT文件分配表的原理得到了正确的理解与应用。

\textbf{任务二}在任务一的基础上引入多进程并发访问机制,通过为每个进程维护独立的进程控制块和文件描述符表,使得多个进程可以同时操作文件系统。针对目录项、FAT表等全局共享资源,使用互斥锁\texttt{fs\_mutex}加以保护;针对普通文件,采用读者-写者问题的经典解法,在读进程和写进程之间实现了正确的同步互斥。测试结果表明:写进程打开文件时其他进程无法访问,多个读进程可以同时读取同一文件,有读进程存在时写进程访问会被拒绝,以上场景均符合预期,验证了并发控制机制的正确性。
% \end{content}
